Los resultados confirman que el tiempo y la memoria de los tres algoritmos crecen según su complejidad teórica, y que cada paradigma ofrece un compromiso distinto entre calidad de la solución y costo computacional. La fuerza bruta es exacta, pero su tiempo exponencial la limita a grillas pequeñas. La programación dinámica también es exacta y reduce el tiempo a un crecimiento lineal en el tamaño de la grilla, a cambio de una memoria proporcional a este. Greedy es el más eficiente, pero su brecha respecto del óptimo aumenta con el tamaño de la grilla y falla en los terrenos con murallas. Por lo tanto, la programación dinámica es la opción que combina optimalidad y escalabilidad, mientras que greedy solo conviene cuando prima la rapidez y se tolera una solucióñ subóptima o inexistente.
